\chapter{绪论}\label{chap:lead}

\section{研究意义与目的}\label{sec:meaning}
机械臂是现代工业自动化及智能制造领域当中不可或缺的关键装备，其应用已经拓展到了精密装配、智慧物流、医疗手术还有家庭服务等许多领域\cite{JDCP202405017, billard2019trends}。我国发布的《“十四五”机器人产业发展规划》\cite{miit2021robot}中明确指出需要重点攻克机器人智能化方面的关键技术，从而推动机器人在国民经济各个领域实现规模化应用。随着人口红利逐步消退以及灵活制造需求的快速增长，工业界对于机械臂自主化以及智能化水平提出了更为严格的要求，既要能够胜任结构化生产线上的高精度操作，也需在非结构化场景中适应各种变化的任务需求\cite{xu2024survey}。但就当前状况来看，机械臂的部署仍在很大程度上依赖于专业技术人员所开展的离线编程以及示教调试工作\cite{alexandrova2014robot, wang2025low}，这种方式部署成本偏高且对调试人员具备专业水平要求，人机交互效率不够理想，制约了机械臂在服务业以及中小制造业等领域的进一步推广。

就当前所采用的机械臂控制方法而言，传统比例积分微分(Proportional-Integral-Derivative, PID)控制、线性二次调节器(Linear Quadratic Regulator, LQR)控制以及模型预测控制(Model Predictive Control, MPC)等方法虽然在结构化环境中表现得比较成熟，但它们对精确动力学建模的依赖让这类方法难以同时兼顾泛化能力以及鲁棒性\cite{elguea2023review}。强化学习(Reinforcement Learning, RL)则借助自主试错与同环境持续交互的方式，能够在一定程度上学到具有鲁棒性的操作策略，不过训练完成之后的策略网络本质上仍然是面向单一任务的数值映射，缺少对人类语义意图的理解，多任务泛化能力偏弱，同时也难以通过自然语言进行直观交互\cite{feng2025embodied}。

\begin{figure}[htbp]
  \centering
  \begin{minipage}[b]{0.45\textwidth}
    \centering
    \includegraphics[width=\textwidth]{figures/robotic arm in factory}
  \end{minipage}
  \quad
  \begin{minipage}[b]{0.40\textwidth}
    \centering
    \includegraphics[width=\textwidth]{figures/robotic arm for daily use.jpg}
  \end{minipage}
  \caption{机械臂在工业生产与家庭服务领域的应用}
  \label{fig:robot-applications}
\end{figure}

近些年来，大语言模型(Large Language Models, LLMs)在自然语言理解、常识推理以及零样本泛化等方面所展现出的能力\cite{zeng2023large}，给机械臂的智能化控制提供了一条新的可行途径。LLMs能够在不需要额外训练的情况下理解用户意图并且开展任务规划，不过它在本质上依然是在离散文本空间中的概率预测，不具备真实世界连续状态空间的直接感知且缺少底层控制方面的知识，同机械臂连续动作的执行之间存在着比较明显的语义鸿沟\cite{bai2024embodied_llm}。如何把LLMs的语义理解能力下延到连续动作空间，是本文所关注的核心问题。

本文旨在探索一种端到端的机械臂通用控制方法，借助LLMs在语义理解以及规划方面的能力，以行为克隆的思路运用监督指令微调把自然语言直接映射为底层动作序列。整个流程主要包含三个阶段：第一步运用强化学习训练各个任务的专家策略，第二步采集高质量轨迹数据并且完成数据校验以及清洗工作，最后凭借参数高效微调让轻量级语言模型学会从文本化观测状态到连续动作的跨模态映射，以此来实现将自然语言作为交互接口的多任务机械臂控制。

\section{国内外研究现状}\label{sec:status}
\subsection{基于深度强化学习的机械臂控制研究进展}\label{sec:rl-survey}
机械臂控制研究已经经历了比较长的发展历程，传统控制方法如PID控制、LQR控制以及MPC在过去很长一段时间内都处于核心地位\cite{borase2021review, chacko2023lqr, schwenzer2021review}。这些方法需要依靠包含机器人动力学以及环境物理特性的精确数学模型来设计控制逻辑，在结构化环境中表现出良好的稳定性以及可解释性。不过，它们的控制性能在很大程度上依赖于建模的精确程度，一旦面对复杂的非线性接触或者动态变化的作业场景，传统方法往往需要进行大量人工调参，自适应能力比较有限\cite{elguea2023review}。

\begin{figure}[htbp]
  \centering
  \includegraphics[width=0.4\textwidth]{figures/rl_workflow}
  \caption{RL算法智能体与环境交互的工作流程示意图}
  \label{fig:rl-workflow}
\end{figure}

为了克服传统方法对精确建模的依赖，RL给机械臂控制提供了一条数据驱动的途径。RL智能体不需要预先知道环境的动力学方程，仅需凭借同环境的持续交互来获取奖励信号，就能够端到端地优化从感知到动作的映射策略\cite{okafor2021heuristic, le2022deep}。这种试错学习机制让机械臂能够自主习得避障以及轨迹平滑等复杂行为，进而降低了它对精确动力学模型的依赖\cite{adams2022survey, milani2024explainable}。

不过，传统RL算法受限于表格型值函数以及线性策略的表征容量，难以处理机械臂控制中高维连续的状态以及动作空间。深度学习的引入拓宽了策略网络的表征能力，让它们能够从原始传感器数据中直接学到复杂的非线性映射\cite{han2023survey, ter2026taxonomy}。2013年，Mnih等\cite{mnih2013playing}提出深度Q网络(Deep Q-Network, DQN)，首次把深度网络同RL进行结合，依靠经验回放以及目标网络机制解决了训练不稳定性的问题，验证了深度网络近似价值函数的可行性。但DQN仅仅适用于离散动作空间，而机械臂操作涉及高维连续关节控制，直接把关节力矩离散化将会导致动作空间维度爆炸以及控制精度的损失\cite{ramamurthy2022reinforcement}，进而催生了面向连续动作空间的策略优化方法研究。

为了能够在连续动作空间中直接优化策略，研究人员先后提出了一系列策略梯度以及Actor-Critic算法。Schulman等\cite{schulman2017proximal}提出的近端策略优化(Proximal Policy Optimization, PPO)算法凭借裁剪目标函数平衡了训练稳定性以及样本效率，已经成为机械臂控制的主要基准算法之一。Lillicrap等\cite{lillicrap2020continuous}提出深度确定性策略梯度(Deep Deterministic Policy Gradient, DDPG)算法，该方法运用演员-评论家架构让神经网络直接输出关节连续力矩，在基于MuJoCo的连续控制仿真环境中取得了比较显著的性能提升。胡春阳等\cite{hu2021rl_visual}在此基础上把单阶段目标检测(You Only Look Once, YOLO)视觉感知同DDPG进行结合，实现了视觉引导下的端到端机械臂自主控制。

针对DDPG在训练过程中出现的Q值过估计问题，Fujimoto等\cite{fujimoto2018addressing}提出双延迟深度确定性(Twin Delayed Deep Deterministic, TD3)算法，引入了双Critic网络以及延迟更新机制，提高了在复杂操作任务当中的鲁棒性\cite{zhang2020td3}。Haarnoja等\cite{haarnoja2018soft}提出软演员-评论家(Soft Actor-Critic, SAC)算法，通过最大化策略熵来增强探索能力，在精细操作中表现出了更高的采样效率以及泛化性能\cite{de2021soft, kaufmann2023survey}。Kuznetsov等\cite{kuznetsov2020tqc}进一步从概率分布的视角提出了截断分位数(Truncated Quantile Critics, TQC)算法，运用多网络分位数回归近似回报分布并且凭借截断机制抑制过估计偏差，在MuJoCo机械臂控制基准上超越了SAC等主流算法。

\begin{figure}[htbp]
  \centering
  \includegraphics[width=\textwidth]{figures/TQC分位数原子截断.png}
  \caption{TQC算法的价值分位数截断机制示意图\cite{kuznetsov2020tqc}}
  \label{fig:tqc-quantile}
\end{figure}

不过，上述RL算法虽然在单任务控制上取得了比较高的成功率，但仍然面临着一些不足。首先是面向特定任务的奖励函数独立训练，多任务泛化能力偏弱；另外策略网络本质上仍然是数值到数值的黑盒映射，没有办法理解自然语言指令以及高级语义意图。出于这些原因，研究者开始尝试把语言理解能力引入机械臂控制，试图打通自然语言同底层动作之间的通道。

\subsection{LLMs在机械臂控制领域的应用现状}\label{sec:llm-survey}
LLMs在海量文本数据上预训练所获得的语义理解以及常识推理能力\cite{zeng2023large}，推动着它们逐步被引入到具身智能的研究当中。传统机械臂控制系统依赖于结构化指令或者预先编程的动作序列，一旦面对开放式的自然语言指令就难以有效应对，而LLMs能够在不进行额外训练的前提下完成意图理解以及任务规划\cite{huang2022language}。近期有不少综述工作从不同视角出发梳理了这一交叉领域的进展，其中Zeng等\cite{zeng2023large}总结了LLMs在机器人系统中的整体应用架构，白辰甲等\cite{bai2024embodied_llm}从感知、推理以及控制三个层面分析了大模型驱动的具身智能融合路径，王文晟等\cite{wang2025embodied_survey}从控制层级以及系统架构维度梳理了基于大模型的具身智能方法，郭庭航等\cite{guo2025llm_robot_planning}则聚焦于LLMs在机器人任务规划中环境感知、行为分解以及运动规划方面的应用。

不过，像GPT-5\cite{singh2025openai}、Qwen3\cite{yang2025qwen3}等主流模型所包含的参数量过大，给机械臂的端侧部署带来了非常高的计算成本以及推理延迟，很难满足高频底层控制对于实时性的要求\cite{ma2025running}。所以不少工作选择把LLMs定义为上层规划器的角色，实现任务层级上的低频调用，在保留高级语义理解能力的同时通过预先的动作封装来实现低延迟的机械臂动作映射。2022年，Ahn等\cite{ahn2022can}提出的SayCan框架把未经微调的540B参数PaLM模型当作高层任务规划器，结合底层技能的价值函数来筛选在当前真实环境中既符合语义逻辑又切实可行的动作，最终在101项操纵任务中取得了大于60\%的成功率。在此基础上，Driess等\cite{driess2023palm}提出的PaLM-E模型额外对原始PaLM模型进行微调，并训练了独立的视觉Transformer(Vision Transformer, ViT)网络用以处理视觉输入，进一步完善了智能体的感知并提升了长时序规划能力，但这两种方法仅对LLMs的上层规划问题进行了优化，底层操作还是依赖预封装的底层控制策略。对于这一问题，Dalal等\cite{dalal2024plan}将运动规划和强化学习融入了底层策略，进一步地解决了长序列控制和规划问题，在包括MetaWorld等25项复杂机械臂控制任务上达到了高于85\%的成功率。

\begin{figure}[htbp]
  \centering
  \includegraphics[width=\textwidth]{figures/plan-seq-learn.png}
  \caption{基于LLMs作为上层策略的控制框架\cite{dalal2024plan}}
  \label{fig:llm-control-framework}
\end{figure}

另外也有不少工作试图把LLMs的代码能力同RL在控制策略探索上的优势进行融合，借助LLMs来进一步引导RL的训练。Xie等\cite{xie2024text2reward}提出的Text2Reward框架驱动GPT-4根据自然语言目标直接编写可执行的Python奖励代码。这项工作在ManiSkill2以及MetaWorld的17项机械臂操纵任务上进行了验证，生成的代码在90\%的情况下能被环境正确转译，成功驱动机器人完成了开门、推椅等任务。Ma等\cite{ma2023eureka}提出的EUREKA算法则把奖励函数设计建模为一个条件搜索问题。该方法把环境代码直接提供给GPT-4，运用大模型的上下文学习能力生成候选奖励函数，并且凭借强化学习的策略表现作为反馈来进行多轮迭代优化。在Isaac Gym环境的29项复杂操纵任务中，EUREKA生成的奖励函数在83\%的任务上超越了人类专家设计的基准。这两项工作说明，运用大参数LLMs的代码先验，可以降低RL中奖励函数设计的难度并且提高采样效率。

针对高质量机械臂真机示教数据获取成本偏高的问题，Ha等\cite{ha2023scaling}提出的Scaling\&Distilling方法展示了LLMs在数据自动化构建方面所具备的潜力。该研究运用LLMs把复杂的机械臂任务拆解为子任务文本，并且在仿真环境中执行以生成成功的运动轨迹，最后把这些轨迹蒸馏到轻量化的动作策略模型当中。基于GPT-3以及LLama2的测试表明，该方法在稀疏奖励任务中的学习效率提高了76\%。

这些工作从任务规划、RL辅助训练以及数据自动化构建等方面推进了LLMs同机械臂控制之间的结合，但它们的共同局限在于依赖分层的转译框架。LLMs仅承担高层语义解析，底层动作执行仍然交由独立并且封闭的控制模块来完成，感知、理解以及执行之间存在着信息瓶颈。所以，研究者开始探索直接运用专家示范数据对模型进行行为克隆式的训练，以专家的状态动作对为监督信号微调整个控制链路，也就催生了视觉语言动作(Vision-Language-Action, VLA)等端到端方法来实现自然语言控制。

\subsection{基于行为克隆的机器人控制策略学习进展}\label{sec:il-survey}
行为克隆(Behavior Cloning, BC)是模仿学习中最基础并且应用最为广泛的方法，它的核心思路在于把策略学习转化为监督学习问题。BC方法虽然实现起来比较简单，但却存在分布偏移的问题，也就是策略的微小误差会在推理阶段不断累积导致轨迹逐渐偏离训练分布\cite{ross2011reduction}。针对这一问题，Ross等\cite{ross2011reduction}提出DAgger算法，通过迭代地在策略自身访问的状态上查询专家标签，有效地缓解了累积误差。这些基础方法为后续运用大规模数据的机器人策略学习提供了理论框架。

\begin{figure}[htbp]
  \centering
  \includegraphics[width=\textwidth]{figures/OpenVLA.png}
  \caption{OpenVLA架构\cite{kimOpenVLAOpenSourceVisionLanguageAction}}
  \label{fig:openvla}
\end{figure}

在数据规模化以及感知输入多元化的趋势下，基于行为克隆的机器人控制方法取得了比较大的突破。通过在大规模专家轨迹数据集上，以观测动作对为监督信号对模型进行端到端微调，让模型在给定感知输入以及语言指令条件下输出正确的动作序列。2022年，谷歌DeepMind提出的RT-1\cite{brohan2022rt}运用EfficientNet提取视觉特征，运用FiLM(Feature-wise Linear Modulation)机制融合语言嵌入，由Transformer输出离散化动作token，在包含13万条真实世界轨迹、覆盖700余项任务的行为克隆数据集上展现出良好的跨任务泛化能力。RT-2\cite{zitkovich2023rt}进一步把大规模预训练视觉语言模型当作基础网络，运用互联网图文数据赋予了机器人更强的常识推理能力，在面对训练集外的新型物体以及复杂语义指令时成功率提高了至少35\%。随着Open X-Embodiment等包含22种机器人平台、超过100万条轨迹的大规模具身数据集的发布\cite{oneillOpenXEmbodimentRobotic2024}，OpenVLA\cite{kimOpenVLAOpenSourceVisionLanguageAction}以7B参数的Llama 2为语义核心，融合DINOv2以及SigLIP视觉特征，经97万条轨迹的行为克隆微调后在多项操纵基准测试中超越了闭源的RT-2-X。Octo模型\cite{octomodelteam2024octoopensourcegeneralistrobot}同样基于行为克隆范式，引入扩散策略头处理多峰动作分布，利用T5编码器来处理指令，解决了自回归模型在多模态动作生成中的局限性。

\begin{figure}[htbp]
  \centering
  \includegraphics[width=\textwidth]{figures/diffusion_summary.png}
  \caption{扩散策略类别与具体架构\cite{li2026diffusion_survey}}
  \label{fig:diffusion-summary}
\end{figure}

在动作表示方面，近期的研究围绕动作序列建模以及生成式方法展开了深入探索。Zhao等\cite{zhao2023learning}提出ACT算法，凭借预测未来多个时间步的动作块而非单步动作，有效地缓解了决策频率以及任务完成效率之间的矛盾，在双臂操纵平台上仅用约50条示范即达到了较高的任务成功率。Chi等\cite{chi2023diffusion}提出扩散策略，运用扩散去噪模型生成连续动作序列，天然适应多峰动作分布，在多个操纵基准测试中达到了领先性能。李牧等\cite{li2026diffusion_survey}系统地综述了扩散模型在机器人行为克隆中的应用，指出扩散策略在处理接触丰富型操作任务时相比确定性策略具备显著优势。在此基础之上，RDT-1B\cite{liu2025rdt1bdiffusionfoundationmodel}运用1B参数规模的DiT架构，把扩散去噪集成到Transformer解码器中实现高维动作序列生成。Physical Intelligence开发的$\pi_0$模型\cite{black2024pi0}基于PaliGemma引入流匹配机制，以50Hz频率输出连续动作块，提高了精细操纵的平滑度。$\pi_{0.5}$\cite{intelligence2025pi05visionlanguageactionmodelopenworld}进一步凭借分层架构把任务分解以及底层动作执行进行结合。

行为克隆方法已经从早期的小规模示范学习发展为融合视觉语言动作的大规模基础模型范式，动作生成的连续性以及实时性得到了比较明显的改善。但现有的VLA模型在端侧部署时计算资源消耗比较大、闭环控制延迟偏高，并且小参数模型在动作向量映射的一致性以及语义跟随精度上尚有不足。关于如何运用高质量的专家数据在有限参数规模下实现轻量模型的有效对齐，目前仍然缺乏成熟的解决方案。

\section{当前通用机械臂控制模型的研究瓶颈}\label{sec:bottleneck}
综合以上研究现状来看，当前通用机械臂控制模型的研发主要面临着以下几个方面瓶颈。

第一，高层语义理解以及底层连续控制之间的语义鸿沟还没有得到有效的弥合。虽然LLMs具备比较强的自然语言理解以及任务规划能力，但其输出仍然是离散文本符号序列，无法直接驱动机械臂的连续关节运动。现有的SayCan、PaLM-E等分层框架虽然凭借转译层把语言指令映射为底层技能调用，但中间转译环节引入了信息瓶颈，最终的性能也需要依赖底层控制策略的能力，限制了端到端的控制精度以及响应速度。

第二，在多任务泛化能力上存在不足。当前基于RL的机械臂控制方法通常面向单一任务训练独立策略，一旦面对新任务就需要重新设计奖励函数并且从头训练，部署成本偏高并且难以进行扩展。而基于VLA的方法虽然具备跨任务泛化能力，但它们依赖大规模的真机视觉轨迹数据，数据获取成本非常高，并且对于未训练过的新任务完成能力也不够。

第三，端侧部署的计算效率以及实时控制需求之间存在着比较突出的矛盾。OpenVLA、$\pi_0$等高性能VLA模型参数规模达到数亿甚至数十亿量级，对部署资源的需求比较高，并且推理延迟很难满足机械臂闭环控制对高频输出的实时性要求。而轻量化模型在动作预测的精度以及一致性方面尚有不足，制约了它们在实际场景当中的部署应用。

针对上述瓶颈，本文提出了一种基于行为克隆以及监督微调的轻量级技术路线。具体做法是先运用强化学习训练高质量的单任务专家策略，以低成本获取标准轨迹数据，再以专家轨迹为监督信号，凭借参数高效微调让轻量级LLM学习从自然语言指令以及数值观测到连续动作序列的跨模态映射，从而在保持模型轻量化的同时实现多任务的统一控制。

\section{主要研究内容及文章结构}\label{sec:structure}
本文围绕基于轻量大语言模型的机械臂多任务通用控制问题，提出了从强化学习专家策略训练到轻量级语言模型指令微调的技术路线。具体来说，本文首先凭借TQC算法以及HER机制训练单任务专家策略以获取高质量轨迹数据，随后设计从物理状态到自然语言的数据转化以及轨迹清洗方法构建多任务微调数据集，最终凭借LoRA参数高效微调让轻量级语言模型学会从自然语言指令到连续动作的跨模态映射，并且凭借vLLM推理引擎实现低延迟闭环推理。本文研究内容共分为五章，如\cref{fig:thesis-structure}所示。

\begin{figure}[htbp]
  \centering
  \includegraphics[width=\textwidth]{figures/workflow}
  \caption{论文整体文章结构与技术路线}
  \label{fig:thesis-structure}
\end{figure}

\cref{chap:lead}探讨了基于自然语言的机械臂多任务通用控制方法的研究背景与意义。章节详细回顾了深度强化学习在机械臂控制中的研究进展，大语言模型在机械臂控制领域的应用现状，以及基于行为克隆的机器人控制策略学习进展。在此基础之上，对当前通用机械臂控制所面临的若干研究瓶颈进行了分析，同时概述了本文的主要研究内容以及章节安排。

\cref{chap:reinforce}研究了基于TQC算法的机械臂操作任务专家策略学习问题。该章首先概述强化学习基础理论，介绍马尔可夫决策过程以及策略优化的基本框架；然后阐述TQC算法与HER机制的原理；最终针对Fetch机械臂的Reach、Push、Slide以及PickAndPlace四项操作任务，分别设计状态空间、动作空间以及奖励函数，训练单任务专家策略。该章节最后凭借消融实验对比TQC以及TQC+HER的训练表现，验证HER在不同任务复杂度下的作用。

\cref{chap:dataset}研究了面向大语言模型微调的多任务数据集构建以及清洗问题。该章设计了从物理状态到自然语言的数据转化机制，通过将采集到的专家策略实际运行的原始轨迹数据经三角色对话模板封装，补充物理量语义解释，从而把专家策略输出的数值轨迹转化为结构化自然语言文本。在此基础之上，提出基于Mujoco仿真引擎闭环回放的轨迹采集以及清洗方法，对转化后的轨迹进行可复现性验证，并且从匹配率、漂移误差以及逐步误差三个维度对清洗后的数据集进行定量质量分析。

\cref{chap:finetune}研究了面向机械臂底层控制策略的轻量LLM指令微调以及闭环推理问题。该章阐述了监督指令微调以及LoRA参数高效微调的基本原理，引入了基于PagedAttention机制的vLLM推理引擎来满足低延迟闭环控制需求，并且把微调以及推理模块整合成为包含观测获取、文本编码、模型推理、动作解析以及环境执行五个阶段的闭环控制流程。该章在6个轻量级基座模型上分别开展单任务以及多任务混合微调实验，从任务成功率以及推理延迟两个维度评估系统性能，分析跨任务数据混合对不同性能水平模型的影响以及vLLM部署的推理加速效果。

\cref{chap:conclusion}对全文研究工作进行了总结，概括了各章节的主要研究成果以及创新点，同时也分析了当前工作在多任务微调策略、视觉感知融合、仿真到真实迁移等方面的局限性，并且对未来研究方向进行了展望。